\documentclass[12pt, a4paper]{article}

\usepackage[a4paper, margin=1in]{geometry}
\usepackage[utf8]{inputenc}
\usepackage[T1]{fontenc}
\usepackage[bahasa]{babel}
\usepackage{amsmath, amssymb}
\usepackage{enumitem}
\usepackage{hyperref}
\usepackage{arabtex}

\title{\textbf{MODUL UTUH PENDIDIKAN AGAMA ISLAM\\ DAN BUDI PEKERTI (PABP) KELAS 10}\\ \large Pendalaman Materi Akademik Berbasis Kisi-Kisi ASAT\\ Kurikulum Merdeka Semester Genap}
\date{}

\begin{document}

\maketitle

\tableofcontents
\newpage

% ==============================================================
\section{BAB I: MENGHINDARI AKHLAK MADZMUMAH DAN MEMBIASAKAN AKHLAK MAHMUDAH}
% ==============================================================

\subsection{Sifat Temperamental (Ghadhab)}

\subsubsection{Pengertian dan Hakikat Ghadhab}
Secara etimologis, \textit{ghadhab} berarti keras, kasar, atau marah. Secara terminologi akademik, \textit{ghadhab} adalah gejolak emosional negatif di dalam hati yang menggerakkan darah untuk meluap, menuntut kepuasan berupa agresi fisik atau verbal, serta melumpuhkan fungsi kontrol akal sehat manusia. Sifat temperamental ini dikategorikan sebagai akhlak madzmumah (tercela) karena menjadi akar dari berbagai tindakan destruktif lainnya, baik bagi diri sendiri maupun tatanan sosial masyarakat.

\subsubsection{Penyebab Sifat Temperamental (Ghadhab)}
Faktor pencetus sifat temperamental dapat dianalisis dari dua dimensi utama:
\begin{enumerate}
    \item \textbf{Faktor Internal (Psikologis):} Adanya penyakit hati seperti kesombongan (\textit{kibar}), merasa diri paling benar atau mulia (\textit{ujub}), egoisme tinggi, serta kurangnya latihan pengendalian diri (\textit{mujahadatun nafs}).
    \item \textbf{Faktor Eksternal (Lingkungan):} Adanya provokasi atau ejekan dari pihak lain, kelelahan fisik yang akumulatif, lingkungan sosial yang terbiasa menyelesaikan konflik dengan kekerasan, serta kondisi lingkungan yang tidak berjalan sesuai dengan ekspektasi pribadi.
\end{enumerate}

\subsubsection{Metode Menghindari Sifat Ghadhab}
Islam memberikan panduan kuratif dan preventif untuk meredam ghadhab agar kembali pada ketenangan jiwa:
\begin{itemize}
    \item \textbf{Membaca Ta'awwudz:} Memohon perlindungan kepada Allah SWT untuk mengusir bisikan setan yang menyalakan api amarah.
    \item \textbf{Mengubah Posisi Tubuh:} Jika emosi memuncak dalam posisi berdiri, hendaknya segera duduk; jika belum mereda, hendaknya mengambil posisi berbaring sesuai tuntunan sunnah.
    \item \textbf{Berwudhu dengan Air Mengalir:} Membasuh tubuh dengan air suci untuk mendinginkan gejolak darah dan emosi yang secara metafisik digerakkan oleh unsur api setan.
    \item \textbf{Menjaga Lisan (Diam):} Menahan diri dari berbicara saat marah guna mencegah keluarnya kata-kata makian atau keputusan yang merugikan.
\end{itemize}

\subsection{Perilaku Berani Membela Kebenaran (Syaja'ah)}

\subsubsection{Pengertian dan Konsep Syaja'ah}
\textit{Syaja'ah} secara bahasa berarti berani, gagah, atau perkasa. Secara istilah syariat, \textit{syaja'ah} adalah kekuatan hati, keteguhan mental, dan keberanian yang proporsional dalam menegakkan kebenaran dan keadilan demi meraih ridha Allah SWT, tanpa terpengaruh oleh rasa takut terhadap makhluk. Dalam konsep etika Islam, \textit{syaja'ah} merupakan titik keseimbangan (\textit{al-wasath}) yang ideal di antara dua kutub ekstrem yang tercela, yaitu sifat pengecut (\textit{jubn}) dan sifat nekat atau membabi buta (\textit{tahawwur}).

\subsubsection{Faktor Pembentuk Sikap Syaja'ah}
Keberanian sejati tidak lahir dari kekuatan fisik semata, melainkan dibentuk oleh fondasi spiritual berikut:
\begin{itemize}
    \item \textbf{Tauhid yang Murni:} Keyakinan mutlak bahwa hidup, mati, rezeki, dan segala mudharat berada sepenuhnya di bawah ketetapan Allah SWT.
    \item \textbf{Orientasi Akhirat:} Mengutamakan balasan surga dan rida Ilahi di atas kenyamanan material duniawi yang fana.
    \item \textbf{Ketenangan Jiwa (\textit{Sakinah}):} Hati yang bersih dari ketergantungan kepada makhluk melahirkan ketabahan dalam menghadapi ujian berat.
    \item \textbf{Kecintaan pada Keadilan:} Rasa tanggung jawab moral yang tinggi untuk menghapuskan kezaliman di muka bumi.
\end{itemize}

\subsubsection{Hikmah, Manfaat, dan Implementasi Syaja'ah dalam Kehidupan}
Penerapan karakter \textit{syaja'ah} memberikan implikasi positif yang sangat luas:
\begin{enumerate}
    \item \textbf{Hikmah Individual:} Membentuk pribadi yang berintegritas, tegar, berwibawa, jujur, serta konsisten dalam menolak segala bentuk kemaksiatan.
    \item \textbf{Manfaat Sosial:} Terwujudnya stabilitas keamanan hukum, terlindunginya kaum yang lemah dari kezaliman struktural, dan tegaknya pilar \textit{amar ma'ruf nahi munkar}.
    \item \textbf{Implementasi Nyata:} Berani mengakui kesalahan diri sendiri, berani menyampaikan kebenaran secara santun kepada penguasa yang zalim (\textit{kalimatu haqqin 'inda sulthanin ja'ir}), serta berani menolak ajakan pergaulan bebas di lingkungan sekolah.
\end{enumerate}

\subsection{Kontrol Diri (Mujahadatun Nafs)}

\subsubsection{Hakikat Mujahadatun Nafs}
\textit{Mujahadatun nafs} diartikan sebagai perjuangan sungguh-sungguh melawan kecenderungan hawa nafsu internal yang buruk (\textit{nafsul ammarah bis-su'}) agar tunduk pada kendali akal sehat dan bimbingan syariat. Kontrol diri ini merupakan poros utama pembentukan karakter mulia (akhlak mahmudah).

\subsubsection{Pembiasaan Perilaku Kontrol Diri}
Langkah konkret untuk membudayakan perilaku kontrol diri dalam keseharian meliputi:
\begin{itemize}
    \item \textbf{Manajemen Emosi Aktual:} Melakukan jeda berpikir sebelum merespons stimulus eksternal agar tidak terjebak dalam penyesalan.
    \item \textbf{Melatih Disiplin Lewat Puasa:} Mengondisikan diri menahan keinginan biologis sah (makan dan minum) demi mengasah ketajaman spiritual.
    \item \textbf{Muhasabah Rutin:} Mengevaluasi niat dari setiap perbuatan yang telah dikerjakan agar senantiasa terjaga keikhlasannya.
\end{itemize}

\newpage

% ==============================================================
\section{BAB II: MENERAPKAN AL-KULLIYATU AL-KHAMSAH DALAM KEHIDUPAN SEHARI-HARI}
% ==============================================================

\subsection{Pengertian al-Kulliyatu al-Khamsah}
\textit{Al-Kulliyatu al-Khamsah} secara bahasa bermakna lima prinsip universal. Dalam kajian Ushul Fikih dan filosofi hukum Islam (\textit{Maqashid asy-Syari'ah}), istilah ini merujuk pada lima kepentingan mutlak (\textit{dharuriyyat}) kemanusiaan yang wajib dijaga, dilindungi, dan dipelihara secara utuh agar manusia dapat hidup dengan terhormat dan selamat di dunia maupun akhirat. Pengabaian terhadap kelima elemen pokok ini akan memicu kerusakan sistemik (\textit{fasad}) dalam kehidupan individu dan tatanan sosial.

\subsection{Macam-Macam al-Kulliyatu al-Khamsah dan Cara Menjaganya}

\subsubsection{1. Menjaga Agama (\textit{Hifzhu al-Din})}
Menjaga agama menempati tangga prioritas tertinggi karena menjadi ruh dari orientasi penciptaan manusia. Perlindungan ini dilakukan melalui dua pendekatan metodis:
\begin{itemize}
    \item \textbf{Aspek Konstruktif (\textit{Min janib al-wujud}):} Menegakkan rukun Islam secara konsisten, memperdalam kajian akidah shahihah, serta menyebarkan syiar dakwah secara damai.
    \item \textbf{Aspek Defensif (\textit{Min janib al-'adam}):} Membentengi diri dari pemikiran menyimpang (sekularisme radikal, liberalisme agama), menjauhi kesyirikan, serta menolak segala bentuk penistaan terhadap syariat.
\end{itemize}

\subsubsection{2. Menjaga Jiwa (\textit{Hifzhu al-Nafs})}
Jiwa manusia memiliki kehormatan tertinggi di hadapan Allah SWT. Penjagaannya diatur melalui instrumen normatif berikut:
\begin{itemize}
    \item Larangan mutlak melakukan pembunuhan, penganiayaan fisik, mutilasi, maupun tindakan bunuh diri.
    \item Pemberlakuan sistem hukum sanksi \textit{qishash} (balasan setimpal) dan \textit{diyat} (denda materi) untuk mencegah tindakan kriminalitas pembunuhan di masyarakat.
    \item Kewajiban menjaga kesehatan tubuh melalui konsumsi makanan bergizi, pemenuhan layanan medis, dan pemeliharaan keselamatan lingkungan kerja.
\end{itemize}

\subsubsection{3. Menjaga Akal (\textit{Hifzhu al-Aql})}
Akal merupakan basis pembebanan syariat (\textit{manathat taklif}) yang membedakan derajat manusia dengan makhluk lainnya. Upaya pemeliharaannya diwujudkan dengan:
\begin{itemize}
    \item Kewajiban menuntut ilmu (\textit{thalabul 'ilmi}) secara kontinu guna mengembangkan kapasitas intelektual dan kognitif.
    \item Pengharaman mutlak atas konsumsi \textit{khamr}, narkotika, psikotropika, zat adiktif, serta aktivitas lain yang merusak susunan saraf pusat otak.
    \item Penetapan sanksi hukum \textit{hadd} bagi penyalahguna zat memabukkan demi memelihara kesadaran berpikir kolektif.
\end{itemize}

\subsubsection{4. Menjaga Keturunan (\textit{Hifzhu al-Nasl})}
Prinsip ini berfokus pada pelestarian eksistensi umat manusia, kesucian garis nasab, dan keluhuran moral keluarga. Caranya meliputi:
\begin{itemize}
    \item Mensyariatkan institusi pernikahan yang sah dan suci sebagai gerbang legal-spiritual interaksi biologis pria dan wanita.
    \item Mengharamkan perilaku pergaulan bebas, perbuatan zina, prostitusi, pornografi, dan penyimpangan seksual yang merusak garis keturunan dan moralitas sosial.
    \item Mewajibkan orang tua memberikan nafkah halal dan pendidikan akhlakul karimah yang komprehensif bagi anak cucunya.
\end{itemize}

\subsubsection{5. Menjaga Harta (\textit{Hifzhu al-Maal})}
Harta diposisikan sebagai sarana penunjang peribadatan dan kesejahteraan hidup. Regulasi Islam untuk melindunginya adalah:
\begin{itemize}
    \item Memotivasi umat untuk bekerja mencari rezeki secara produktif, halal, jujur, dan berlandaskan profesionalisme.
    \item Mengharamkan praktik pencurian, perampokan, korupsi, manipulasi data keuangan, perjudian, serta transaksi berbasis riba.
    \item Menerapkan sanksi hukum yang tegas bagi pelaku kejahatan ekonomi guna menjamin hak kepemilikan aset privat dan publik.
\end{itemize}

\newpage

% ==============================================================
\section{BAB III: PERAN TOKOH ULAMA DALAM PENYEBARAN ISLAM DI INDONESIA (DAKWAH WALI SONGO)}
% ==============================================================

\subsection{Strategi Dakwah Wali Songo di Tanah Jawa}
Proses Islamisasi di Nusantara, khususnya di tanah Jawa yang diarsiteki oleh Dewan Wali Songo, merupakan potret keberhasilan dakwah kultural terbaik di dunia. Mereka tidak menggunakan jalur militer atau kekerasan, melainkan menerapkan metode \textit{dakwah bil-hikmah} (kebijaksanaan) and \textit{mau'izhatul hasanah} (nasihat yang baik). Strategi utamanya berfokus pada pribumisasi Islam—mengintegrasikan nilai-nilai tauhid ke dalam kebudayaan masyarakat Nusantara tanpa merusak struktur tradisi yang telah mengakar kuat sebelumnya.

\subsection{Metode Dakwah Spesifik Tokoh Wali Songo}

\begin{itemize}[leftmargin=*]
    \item \textbf{Sunan Gresik (Maulana Malik Ibrahim):} Dianggap sebagai pelopor Wali Songo. Beliau menggunakan strategi pendekatan ekonomi melalui jalur perdagangan internasional dan pembaruan sistem pertanian rakyat. Beliau juga membuka panti pengobatan gratis dan merangkul kasta Sudra untuk menghapuskan diskriminasi sosial lewat prinsip kesetaraan dalam Islam.
    
    \item \textbf{Sunan Ampel (Raden Rahmat):} Perancang \textit{blueprint} kaderisasi dai di tanah Jawa dan pendiri Pesantren Ampeldenta. Beliau sangat terkenal dengan metode doktrin moral \textbf{"Moh Limo"} (Lima Pantangan): \textit{Moh Main} (tidak berjudi), \textit{Moh Ngombe} (tidak mabuk), \textit{Moh Maling} (tidak mencuri), \textit{Moh Madat} (tidak mengonsumsi narkoba), dan \textit{Moh Madon} (tidak berzina).
    
    \item \textbf{Sunan Drajat (Raden Qasim):} Pendekatan dakwah beliau bertumpu pada filantropi Islam, pemberdayaan ekonomi, dan jaminan sosial bagi kaum miskin serta anak yatim. Beliau terkenal dengan pitutur luhurnya: memberikan payung kepada yang kehujanan, memberikan pakaian kepada yang telanjang, dan memberikan makan kepada yang lapar. Beliau juga menggagas media tembang \textit{Pangkur}.
    
    \item \textbf{Sunan Kudus (Ja'far Shadiq):} Memiliki keahlian tinggi dalam bidang fikih dan strategi pertahanan. Pendekatan dakwah beliau sangat menghormati psikologi kultural penganut Hindu-Buddha. Beliau melarang penyembelihan sapi di wilayah Kudus sebagai bentuk toleransi kultural, serta mendesain arsitektur Menara Masjid Kudus menyerupai candi Hindu.
    
    \item \textbf{Sunan Giri (Raden Paku):} Mengembangkan jalur dakwah berbasis institusi pendidikan pesantren dan kekuatan politik struktural (Kerajaan Giri Kedaton). Jaringan santrinya mampu menyebarkan Islam hingga ke wilayah Maluku, Ternate, dan Tidore. Beliau juga menciptakan permainan edukatif anak Islami seperti \textit{Jelungan} dan tembang \textit{Asmarandana}.
    
    \item \textbf{Sunan Kalijaga (Raden Mas Said):} Tokoh seniman dan budayawan agung Wali Songo. Beliau melakukan modifikasi mendalam terhadap kesenian wayang kulit, instrumen gamelan, dan seni ukir agar bermuatan falsafah tauhid. Beliau adalah pencipta tembang mistikus-edukatif yang legendaris, \textit{Ilir-Ilir} dan \textit{Gundul-Gundul Pacul}.
    
    \item \textbf{Sunan Gunung Jati (Syarif Hidayatullah):} Satu-satunya Wali Songo yang memegang kendali pemerintahan secara struktural sebagai Sultan Cirebon. Beliau mengombinasikan misi dakwah dengan perluasan wilayah, penguatan armada militer maritim, diplomasi internasional, serta penyebaran Islam di wilayah Sunda (Jawa Barat) melalui pendekatan sosial-politik.
\end{itemize}

\newpage

% ==============================================================
\section{BAB IV: MENJAUHI PERGAULAN BEBAS DAN PERBUATAN ZINA}
% ==============================================================

\subsection{Kajian Q.S. an-Nur/24: 2}

\subsubsection{Teks Arab dan Transliterasi Sistem}
\begin{center}
\begin{arabtext}
al-zAniyatu wa-al-zAnI fa-ajlidU kulla wA_hidin min-humA mi'ata jaldatin wa-lA ta'_khu_dkum bi-himA ra'fatun fI dIni AllAhi in kuntum tu'minUna bi-AllAhi wa-al-yawmi al-A_khiri wa-l-ya\$had `a_dAbahumA tA'ifatun mina al-mu'minIna.
\end{arabtext}
\end{center}

\subsubsection{Terjemahan Maknawi}
"Pezina perempuan dan pezina laki-laki, deralah masing-masing dari keduanya seratus kali, dan janganlah rasa belas kasihan kepada keduanya mencegah kamu untuk (menjalankan) agama (hukum) Allah, jika kamu beriman kepada Allah dan hari kemudian; dan hendaklah (pelaksanaan) hukuman mereka disaksikan oleh sebagian orang-orang yang beriman."

\subsection{Menelaah Tafsir dan Bentuk Pergaulan Bebas Kontemporer}
Merujuk pada Tafsir Al-Misbah, ayat ini memberikan ketegasan hukum yang mutlak demi memelihara kesucian institusi keluarga dan menjaga martabat kemanusiaan. Hukum cambuk (dera) 100 kali dijatuhkan bagi pelaku zina \textit{ghairu muhshan} (belum pernah terikat pernikahan sah). Perintah agar eksekusi hukuman disaksikan oleh sekelompok orang beriman bertujuan memberikan efek jera secara psikologis dan sosial agar masyarakat tidak menganggap enteng dosa seksual.

Di era modern, pergaulan bebas berkembang melalui berbagai manifestasi yang permisif:
\begin{itemize}
    \item \textbf{Hubungan Pacaran Tanpa Batas Moral:} Aktivitas kedekatan fisik di luar pernikahan yang mengadopsi gaya hidup sekular-hedonis.
    \item \textbf{Seks Bebas (\textit{Free Sex}):} Melakukan hubungan seksual pranikah, baik atas motif suka sama suka maupun komersial.
    \item \textbf{Gaya Hidup Kehidupan Malam (\textit{Nightlife}):} Menghabiskan waktu di tempat hiburan malam yang memfasilitasi konsumsi alkohol dan narkotika bersamaan dengan interaksi bebas pria dan wanita.
\end{itemize}

\subsection{Dampak Destruktif Perbuatan Zina}
Zina dalam Al-Qur'an dilabeli sebagai \textit{fahisyah} (perbuatan keji dan buruk). Dampak buruknya meliputi multi-dimensi kehidupan:
\begin{enumerate}
    \item \textbf{Dampak Kesehatan Medis:} Menjadi vektor utama penyebaran Penyakit Menular Seksual (PMS) yang mematikan, seperti HIV/AIDS, sifilis, kemandulan, dan kanker serviks.
    \item \textbf{Dampak Psikologis:} Menimbulkan perasaan bersalah menahun, depresi akut, degradasi harga diri, serta memicu tindakan aborsi ilegal yang membahayakan nyawa.
    \item \textbf{Dampak Sosial:} Merusak struktur nasab hukum anak, kehancuran keharmonisan keluarga besar, serta mendatangkan laknat, kemiskinan sistemik, dan murka Allah SWT atas suatu negeri.
\end{enumerate}

\newpage

% ==============================================================
\section{BAB V: HAKIKAT MENCINTAI ALLAH SWT., KHAUF, RAJA', DAN TAWAKKAL}
% ==============================================================

\subsection{Hakikat Mencintai Allah SWT (\textit{Mahabbatullah})}
\textit{Mahabbatullah} merupakan kedudukan spiritual (\textit{maqam}) tertinggi dalam pengabdian seorang hamba. Mencintai Allah berarti mengarahkan seluruh kecenderungan hati, raga, dan kehendak pribadi untuk tunduk dan patuh sepenuhnya kepada perintah Allah atas dasar rasa kagum, rindu, dan syukur yang tiada tara.

\subsubsection{Tanda-Tanda Cinta kepada Allah SWT}
\begin{itemize}
    \item Selalu mengutamakan hal-hal yang disukai Allah di atas syahwat atau kesenangan pribadi egoistis.
    \item Merasakan kelezatan rohani, ketenangan (\textit{tumaninah}), dan kekhusyukan yang dalam saat menegakkan salat, berzikir, dan membaca Al-Qur'an.
    \item Mencintai Rasulullah SAW, mencintai sesama mukmin, dan menyayangi makhluk hidup lainnya semata-mata karena Allah.
    \item Rela mengorbankan waktu, harta benda, bahkan jiwa demi memperjuangkan nilai-nilai kebenaran Islam.
\end{itemize}

\subsubsection{Cara Meningkatkan Cinta kepada Allah SWT}
\begin{enumerate}
    \item \textbf{\textit{Tafakur fi al-Khalq}:} Merenungkan keindahan, kompleksitas, dan keagungan penciptaan alam semesta sebagai cerminan sifat-sifat Allah yang Maha Agung.
    \item \textbf{\textit{Tadakkur an-Ni'am}:} Menyadari secara mendalam bahwa setiap helaan napas, kesehatan fisik, kecerdasan akal, dan hidayah iman adalah karunia murni dari Allah yang tak ternilai harganya.
    \item \textbf{Istiqamah Menjalankan Amalan Sunnah:} Mengiringi ibadah wajib dengan salat malam (\textit{tahajjud}), puasa sunnah, serta memperbanyak lantunan \textit{Asmaul Husna}.
\end{enumerate}

\subsection{Hakikat Takut Kepada Allah SWT (\textit{Khauf})}
\textit{Khauf} adalah kegelisahan internal hati atau rasa takut yang mendalam terhadap siksaan, keadilan hukum, dan kemurkaan Allah akibat tumpukan dosa-dosa yang telah diperbuat. \textit{Khauf} yang benar bukan membuat seorang hamba menjauh, melainkan justru membuatnya berlari mendekat demi mencari perlindungan di balik tirai ampunan Allah.
\begin{itemize}
    \item \textbf{Tanda-Tanda Khauf:} Tercegahnya lisan dari ghibah dan dusta, terjaganya pandangan dari maksiat visual, serta segera melakukan tobat nasuha ketika terlanjur melakukan kesalahan.
\end{itemize}

\subsection{Hakikat Berharap Kepada Allah SWT (\textit{Raja'})}
\textit{Raja'} adalah optimisme hati yang dipenuhi harapan besar akan luasnya rahmat, ampunan, dan balasan surga Allah SWT. Sifat \textit{raja'} harus selalu berpasangan dengan amal saleh; harapan tanpa adanya usaha nyata dinilai sebagai angan-angan kosong (\textit{tamanni}) yang tercela. Integrasi antara \textit{khauf} dan \textit{raja'} diibaratkan seperti dua sayap burung; keduanya harus seimbang agar perjalanan spiritual hamba menuju Allah tidak jatuh.

\subsection{Hakikat dan Manfaat Tawakkal}
\textit{Tawakkal} adalah penyerahan total hasil akhir dari segala urusan kepada Allah SWT setelah melakukan ikhtiar atau usaha secara maksimal, cerdas, terencana, dan halal. Tawakkal merupakan bukti nyata keimanan yang matang.
\begin{itemize}
    \item \textbf{Manfaat Tawakkal:} Memberikan ketenangan batin yang luar biasa, menghindarkan diri dari stres atau depresi akut saat menghadapi kegagalan, serta membentengi diri dari sifat sombong (\textit{takabur}) ketika meraih keberhasilan karena menyadari semua itu atas taufik dari Allah SWT.
\end{itemize}

\end{document}